% Zone „kennst du schon“ – Quadratische Gleichungen, Kl. 9 Gymnasium
\setcounter{aufgabe}{0}
\hypertarget{zone}{}\noindent{\large\bfseries Das kennst du schon}\par\medskip

\begin{aufgabe}{Ich kann Zahlen quadrieren, auch negative Zahlen und Dezimalzahlen. Berechne.}
\begin{teilezwei}
\tz $5^2 = $ \leerfeld & \tz $9^2 = $ \leerfeld \\
\tz $(-2{,}5)^2 = $ \leerfeld & \tz $(-7)^2 = $ \leerfeld \\
\tz $(-3)^2 = $ \leerfeld & \tz $-3^2 = $ \leerfeld \\
\tz $2\cdot(-4)^2 = $ \leerfeld & \\
\end{teilezwei}
\end{aufgabe}

\begin{aufgabe}{Ich kann Klammer zuerst und Punkt vor Strich rechnen, auch mit negativen Zahlen. Berechne.}
\begin{teilezwei}
\tz $4\cdot(1 + 2) = $ \leerfeld & \tz $3 + 2\cdot 5 = $ \leerfeld \\
\tz $(-4)\cdot(-4 + 9) = $ \leerfeld & \tz $(-3)\cdot(-3 - 5) = $ \leerfeld \\
\tz $(-6)^2 + 5\cdot(-6) = $ \leerfeld & \tz $2\cdot(-1)^2 - 3\cdot(-1) - 5 = $ \leerfeld \\
\end{teilezwei}
\end{aufgabe}

\begin{aufgabe}{Ich kann eine lineare Gleichung mit Äquivalenzumformungen lösen. Löse die Gleichung. Schreibe die Umformung hinter den Strich.}
\begin{gleichungsraster}[2]
\gl{x + 4 = 9} & \gl{3x = 21} \\ \noalign{\vspace{6pt}}
\gl{x - 2{,}5 = -4} & \gl{-x = 6} \\ \noalign{\vspace{6pt}}
\gl[3]{5x + 3 = 2x - 9} & \gl{x + 7 = 0} \\ \noalign{\vspace{6pt}}
\end{gleichungsraster}
\end{aufgabe}

\begin{aufgabe}{Ich kann durch Einsetzen prüfen, ob eine Zahl eine Gleichung löst. Setze die Zahl ein, rechne beide Seiten aus und kreuze an: wahre Aussage (wA) oder falsche Aussage (fA).}
\begin{teile}
\teil $x = 3$ \ in \ $2x + 1 = 7$ \hfill \kreuz{(wA)}\kreuz{(fA)}
\teil $x = 4$ \ in \ $x - 1 = 5$ \hfill \kreuz{(wA)}\kreuz{(fA)}
\teil $x = -2$ \ in \ $3x + 10 = 4$ \hfill \kreuz{(wA)}\kreuz{(fA)}
\teil $x = -5$ \ in \ $-2x - 3 = 7$ \hfill \kreuz{(wA)}\kreuz{(fA)}
\teil $x = 0{,}5$ \ in \ $4x - 1 = 3$ \hfill \kreuz{(wA)}\kreuz{(fA)}
\end{teile}
\end{aufgabe}

\begin{aufgabe}{Ich kann Quadratwurzeln ziehen und Näherungswerte runden. Ziehe die Wurzel. Geht sie nicht auf, runde auf zwei Nachkommastellen. Gibt es keine Wurzel, schreibe „gibt es nicht“.}
\begin{teilezwei}
\tz $\sqrt{25} = $ \leerfeld & \tz $\sqrt{81} = $ \leerfeld \\
\tz $\sqrt{0{,}49} = $ \leerfeld & \tz $\sqrt{19} \approx $ \leerfeld \\
\tz $\sqrt{-4} = $ \leerfeld & \tz $\sqrt{0} = $ \leerfeld \\
\tz $\sqrt{25 + 144} = $ \leerfeld & \\
\end{teilezwei}
\end{aufgabe}

\begin{aufgabe}{Ich kann an der Scheitelpunktform den Scheitelpunkt und die Öffnung einer Parabel ablesen. Gib den Scheitelpunkt an und kreuze an, wie die Parabel geöffnet ist.}
\begin{teile}
\teil $y = (x - 1)^2 + 2$ \hfill \punktfeld \quad \kreuz{nach oben}\kreuz{nach unten}
\teil $y = (x + 3)^2 - 1$ \hfill \punktfeld \quad \kreuz{nach oben}\kreuz{nach unten}
\teil $y = -(x - 2)^2 + 5$ \hfill \punktfeld \quad \kreuz{nach oben}\kreuz{nach unten}
\teil $y = x^2 - 4$ \hfill \punktfeld \quad \kreuz{nach oben}\kreuz{nach unten}
\end{teile}
Wie viele Nullstellen hat die Parabel? Entscheide am Scheitelpunkt und an der Öffnung, ohne zu rechnen.
\begin{teile}
\teil die Parabel aus a) \hfill \kreuz{keine}\kreuz{eine}\kreuz{zwei}
\teil die Parabel aus c) \hfill \kreuz{keine}\kreuz{eine}\kreuz{zwei}
\teil $y = (x - 4)^2$ \hfill \kreuz{keine}\kreuz{eine}\kreuz{zwei}
\end{teile}
\end{aufgabe}

\begin{aufgabe}{Ich kann ausmultiplizieren, ausklammern und binomische Formeln anwenden. Multipliziere aus und fasse zusammen.}
\begin{gleichungsraster}[1]
\gl{x\cdot(x + 2)} & \gl{-x\cdot(x - 4)} \\ \noalign{\vspace{6pt}}
\gl{(x + 1)^2} & \gl{(x - 3)^2} \\ \noalign{\vspace{6pt}}
\gl{(x + 5)\cdot(x - 2)} & \\ \noalign{\vspace{6pt}}
\end{gleichungsraster}
\medskip
Klammere so viel wie möglich aus.
\begin{gleichungsraster}[1]
\gl{x^2 + 6x} & \gl{3x^2 - 12x} \\ \noalign{\vspace{6pt}}
\end{gleichungsraster}
\end{aufgabe}

\begin{aufgabe}{Ich kann Terme ordnen und zusammenfassen. Fasse zusammen und ordne: zuerst $x^2$, dann $x$, dann die Zahl.}
\begin{gleichungsraster}[1]
\gl{3 + x^2 + 2x} & \gl{5x - 1 + x^2} \\ \noalign{\vspace{6pt}}
\gl{x^2 + 4x - 7 + 2x + 9} & \gl{7 - x^2 + 3x} \\ \noalign{\vspace{6pt}}
\gl{2x^2 - x + 4 - x^2 - 5x + 1} & \\ \noalign{\vspace{6pt}}
\end{gleichungsraster}
\end{aufgabe}

\begin{aufgabe}{Ich finde den Fehler beim Quadrieren negativer Zahlen. Tom rechnet so:}
\rechnung{(-8)^2 + 3\cdot(-8) &= -64 - 24 \\ &= -88}
\begin{teile}
\teil Finde den Fehler, benenne ihn und korrigiere die Rechnung.
\teil Berechne selbst: \ $(-5)^2 + 4\cdot(-5)$
\end{teile}
\end{aufgabe}
